\documentclass[12pt]{article}
\usepackage[utf8]{inputenc}
\usepackage[T1]{fontenc}
\usepackage[margin=1in]{geometry}
\usepackage{setspace}
\usepackage{xcolor}
\usepackage[colorlinks=true, urlcolor=blue!60!black, linkcolor=blue!60!black]{hyperref}
\setlength{\parindent}{2em}
\setlength{\parskip}{1.25em}
\renewcommand{\baselinestretch}{1.0}
\begin{document}
\begin{center}
{\LARGE \textit{Trump v. CASA, Inc.}, 606 U.S. 831 (2025)\par}
\vspace{0.8em}
{\large Case Brief --- Four-Page Edition\par}
\end{center}
\vspace{1em}

\section{Memory Jogger}

Challengers to the 2025 birthright-citizenship executive order won nationwide injunctions in three district courts; without reaching the order's legality, the Court held 6--3 that universal injunctions likely exceed federal courts' equitable authority under the Judiciary Act of 1789, limiting relief to what is needed to give each plaintiff with standing complete relief (\href{https://supreme.justia.com/cases/federal/us/606/24a884/}{Justia}).

\section{Detailed Case Facts}

On January 20, 2025, President Trump issued Executive Order No.~14160, declaring that federal agencies would no longer issue or recognize citizenship documents for children born in the United States whose mother was unlawfully or only temporarily present and whose father was neither a citizen nor a lawful permanent resident. The order gave agencies 30 days to prepare implementation guidance.

Individuals (including pregnant mothers), organizations for immigrants such as CASA, 22 States, the District of Columbia, and San Francisco filed three suits in Maryland, the Western District of Washington, and Massachusetts, alleging violations of the Citizenship Clause and the Nationality Act of 1940.

\section{Procedural History}

All three district courts found the order likely unlawful and entered universal preliminary injunctions barring its enforcement against anyone anywhere. The Fourth, Ninth, and First Circuits denied stays. The Government sought only \textit{partial} stays in the Supreme Court, not defending the order's constitutionality but arguing that the courts lacked authority to protect nonparties. The Court consolidated the applications, heard argument on May 15, 2025, and ruled on June 27, 2025.

\section{Judicial Votes}

Partial stays granted, \textbf{6--3}.
\begin{itemize}
\item \textbf{Majority:} Justice Barrett, joined by Chief Justice Roberts and Justices Thomas, Alito, Gorsuch, and Kavanaugh.
\item \textbf{Concurring:} Justice Thomas (joined by Gorsuch); Justice Alito (joined by Thomas); Justice Kavanaugh.
\item \textbf{Dissenting:} Justice Sotomayor (joined by Kagan and Jackson); Justice Jackson.
\end{itemize}

\section{Holding}

Universal injunctions likely exceed the equitable authority Congress conferred in the Judiciary Act of 1789. The injunctions were stayed ``only to the extent that the injunctions are broader than necessary to provide complete relief to each plaintiff with standing to sue,'' and the order's main provision could not take effect for 30 days. The Court did not decide whether the order is lawful, whether Article III independently bars universal injunctions, or whether the APA authorizes universal vacatur.

\section{Analysis of Opinions}

\textbf{Majority (Barrett).} Section 11 of the Judiciary Act grants jurisdiction over ``all suits \ldots\ in equity,'' and under \textit{Grupo Mexicano} that authority is confined to remedies ``traditionally accorded by courts of equity'' at the founding. Equitable decrees were party-specific, and universal injunctions were ``conspicuously nonexistent'' until the mid-twentieth century. The bill of peace bound small, cohesive groups symmetrically and evolved into the Rule 23 class action; universal injunctions let plaintiffs evade Rule 23 and create one-way litigation in which plaintiffs need win only once. A court may give relief that incidentally benefits nonparties when indivisible relief is needed, as with a neighbor's nuisance, but barring enforcement against one mother's child gives her complete relief. For the States the inquiry is ``more complicated,'' so the Court remanded. An overbroad injunction against a coordinate branch inflicts irreparable harm under \textit{Nken v. Holder}.

\textbf{Concurrences.} Thomas called complete relief a ``ceiling,'' not a ``mandate.'' Alito warned of loopholes in lax state third-party standing and hasty nationwide class certification. Kavanaugh noted that classwide relief and APA set-asides remain available and that the Court itself will decide interim national status.

\textbf{Dissents.} Sotomayor condemned the Government's ``gamesmanship'' in defending a patently unconstitutional order under \textit{Wong Kim Ark}, invoked historical nonparty relief, and urged parents to file class actions. Jackson warned of a ``zone of lawlessness''; Barrett answered that her view embraced an ``imperial Judiciary.''

\textbf{Later.} In \textit{Trump v. Barbara} (2026), the Court affirmed a classwide injunction and held the order violates the Citizenship Clause.

\section{Examples: Future Applications}

\begin{enumerate}
\item \textbf{Same-Side (Stayed): Wire-Transfer Surcharge.} Two plaintiffs, without a class or APA claim, win a nationwide injunction against a surcharge; relief must be limited to their own transfers.
\item \textbf{Same-Side (Narrowed): Contractor Speech Condition.} One contractor's injury is fully redressed by an order protecting it; other contractors may sue or seek a class.
\item \textbf{Opposite-Side (Open Paths): APA Vacatur or a Certified Class.} A trade association's \S~706(2) suit or a properly certified Rule 23(b)(2) class is not barred by \textit{CASA}.
\item \textbf{Opposite-Side (Indivisible): Firing-Range Noise.} Enjoining nighttime firing benefits all neighbors, but there is no way to stop the noise at one house alone.
\item \textbf{Fence-Sitter: State Student-Aid Suit.} A State argues that cross-border student migration makes in-state-only relief unworkable; \textit{CASA} called the state inquiry ``more complicated,'' while Thomas and Alito urged limits.
\end{enumerate}

\section{Critique}

\textbf{Frozen equity.} Critics ask whether freezing federal equity at 1789 Chancery practice makes sense when modern administrative governance and officer suits under \textit{Ex parte Young} had no founding-era twin.

\textbf{Workarounds.} Because the ruling rests on the Judiciary Act rather than Article III, Congress may change it, and plaintiffs have largely replicated nationwide relief through Rule 23(b)(2) classes, state complete-relief showings, and APA vacatur---so \textit{CASA} may simply have shifted the friction to the emergency docket.

\section{Key Quotations}

\begin{quote}
``\,`Complete relief' is not synonymous with `universal relief.' It is a narrower concept: The equitable tradition has long embraced the rule that courts generally `may administer complete relief between the parties.'\,''
\end{quote}
\noindent --- Justice Barrett, for the Court.

\begin{quote}
``But federal courts do not exercise general oversight of the Executive Branch; they resolve cases and controversies consistent with the authority Congress has given them. When a court concludes that the Executive Branch has acted unlawfully, the answer is not for the court to exceed its power, too.''
\end{quote}
\noindent --- Justice Barrett, for the Court.

\begin{quote}
``No right is safe in the new legal regime the Court creates. ... That holding renders constitutional guarantees meaningful in name only for any individuals who are not parties to a lawsuit.''
\end{quote}
\noindent --- Justice Sotomayor, dissenting.

\end{document}
